\chapter{Kajian Pustaka}

\section{Penelitian Terkait}

\section{Estimasi Kedalaman}
Estimasi kedalaman adalah salah satu bagian dari visi komputer yang bertujuan untuk memperkirakan jarak dari objek menggunakan satu atau lebih gambar~\parencite{masoumian2022monocular}.
Kedalaman dapat diestimasi menggunakan beberapa pendekatan, yaitu estimasi kedalaman gambar tunggal (\textit{monocular depth estimation}), estimasi kedalaman binocular, dan estimasi kedalaman \textit{multi-view}~\parencite{masoumian2022monocular}.

\textit{Monocular depth estimation} (MDE) melakukan estimasi kedalaman hanya dengan menggunakan satu gambar RGB, mengurangi kebutuhan data \textit{multi-view} atau perangkat keras khusus seperti kamera stereo.
Hal ini membuat desain sistem yang lebih sederhana dan lebih murah, sehingga lebih ideal pada skenario dunia nyata dibandingkan dengan pendekatan lainnya~\parencite{zhang2025survey}.

Meskipun MDE memiliki keunngulan tersebut, MDE masih merupakan masalah yang sangat menantang, hal ini karena MDE adalah kategori masalah salah-tata (\textit{ill-posed problem}), di mana tidak ada informasi yang dapat diandalkan untuk merekonstruksi kedalaman dari gambar tunggal~\parencite{khan2020deep}.
Hal ini akan menyebabkan hasil MDE akan menghasilkan banyak kemungkinan rekonstruksi 3D.

\section{\textit{Deep Learning} untuk Estimasi Kedalaman}

\section{\textit{Generative Adversarial Networks} (GAN)}

\section{\textit{Conditional GAN} (cGAN) dan \textit{Pix2Pix}}

\section{KITTI \textit{Dataset}}

\section{Fungsi Kerugian (\textit{Loss Function}) pada Estimasi Kedalaman)}

\section{Metrik Evaluasi Kinerja}